\subsection{Constitutive composition and evaluation of the gravitational section}
\label{g26:subsec:evaluacion-viii-en}

The evaluation reuses the generated regional coordinates and record
\(K\) selected in the preceding kernel. Coordinate \(\alpha\)
determines the action sections in
\eqref{v:ant:eq:el-secciones-accion}; in the return section write
\(\eta=\eta_{\rm ret}\) and \(S_*=10^{-34}\,\mathrm{J\,s}\), so
that \(\hbar_{\rm ret}=S_*\eta\). The previously constructed angle
and counterangle give, in the degree representation,
\[
 A_{\deg}=1000\alpha,\qquad C^*_{\deg}=2(\eta+\alpha),\qquad
 x=\frac{\pi_{\rm HMT}A_{\deg}}{180},\quad
 y=\frac{\pi_{\rm HMT}C^*_{\deg}}{180}.
\]
Thus \(\eta=(C^*_{\deg}-2\alpha)/2\). The difference is taken
between coordinates in the same angular representation.

On the generated domain \(0<y<x\), the labelled channels are
\[
 q_\pm=e^{-(x\pm y)},\qquad
 r_\pm=\frac{1-q_\pm^{90}}{1-q_\pm^{120}},\qquad
 0<q_\pm<1.
\]
The quotient comes from the two extension sectors: for
\(S_n(q)=\sum_{j=0}^{n-1}q^j\), the identity
\((1-q)S_n(q)=1-q^n\) proves
\(r_\pm=S_{90}(q_\pm)/S_{120}(q_\pm)\).
The squared channels determine the constitutive readings
\[
 \widehat\varepsilon=r_+^2,\qquad
 \widehat\mu=r_-^2,\qquad
 \widehat c=\frac1{r_+r_-},\qquad
 \widehat Z=\frac{r_-}{r_+}.
\]
For positive velocity, action, and charge bases \(c_*\), \(u_S\),
and \(u_Q\), the dimensional sections are
\[
 \mu=\widehat\mu\frac{u_S}{c_*u_Q^2},\qquad
 \varepsilon=\widehat\varepsilon\frac{u_Q^2}{u_Sc_*},\qquad
 c=c_*\widehat c=\frac1{\sqrt{\mu\varepsilon}}.
\]
Indeed, the product of the first two equations cancels the action and
charge bases, leaving \(\mu\varepsilon=(r_+r_-)^2/c_*^2\); its
positive square root proves the last equality. The factor \(\widehat c\)
transforms the base \(c_*\) once. The SI evaluation below represents
the same \(c\) by \(299792458\,\mathrm{m\,s^{-1}}\) and restores
\(c_*=c/\widehat c\). This operation verifies equivalence of the
dimensional and constitutive representations.

With \(r_N=5759/23040\) and the stated reference \(L_*=1\,\mathrm m\),
substitution of these equations into the radial coefficient gives
\begin{align}
 G&=\frac{c^3L_*^2}{S_*}\frac{r_N^2\alpha^{32}}{\eta}
   =\frac{2c^3L_*^2r_N^2\alpha^{32}}
           {S_*(C^*_{\deg}-2\alpha)},\label{g26:eq:g-angular-viii-en}\\
  &=\frac{L_*^2r_N^2\alpha^{32}}
           {S_*\eta(\mu\varepsilon)^{3/2}}.
 \label{g26:eq:g-constitutivo-viii-en}
\end{align}
These express the same coefficient in the radial realization. The
longitudinal reference \(L_*\), constructed length \(L_\alpha\),
chronological edge \(ct_0\), and horizon radius retain their distinct
types.

Evaluation with the archived regional expansions yields
\begin{align*}
 L_\alpha&=1.6162549826251263301682625013\ldots\,10^{-35}\,\mathrm m,\\
 \hbar_{\rm ret}&=1.0545718169150390611336905847\ldots\,10^{-34}\,
                    \mathrm{J\,s},\\
 t_0&=3.1364999625365490734\ldots\,10^{-45}\,\mathrm s,\\
 G&=6.6742996594140227063677761891\ldots\,10^{-11}\,
                    \mathrm{m^3\,kg^{-1}\,s^{-2}}.
\end{align*}
Calculations at 100 and 140 working digits preserve 71 significant
digits of the evaluation. This numerical stability is distinguished
from experimental uncertainty. Against the CODATA 2022 reference
\(G_{\rm ref}=6.67430(15)\,10^{-11}\,\mathrm{m^3\,kg^{-1}\,s^{-2}}\),
the difference is
\(G-G_{\rm ref}=-3.4058597729\ldots\,10^{-18}\,
\mathrm{m^3\,kg^{-1}\,s^{-2}}\), approximately
\(-0.00227057\) times its standard uncertainty. The reference is
compared after evaluation; rules R4 and R5 and identification R8
retain the scope proved in the preceding section.

The reproduction sources include the regional expansions, record \(K\),
action and counterangle equations, exact verifiers, and decimal evaluation.
They run from the local directory
\path{reproduccion_gravitatoria_20260926} through
\path{verificar_reproduccion.py}. Each test states its domain; reuse of
archived coordinates is distinguished from a new execution of the entire
generator.
